\makeabstract
{
Using Calcium Imaging to Investigate How Flow Direction Influences Habituation in Zebrafish Lateral Line Hair Cells
}
{
Jiahuai Kang (1), Stacey Beganny (1), Allison Platt (1), Josef G. Trapani (1)
}
{
\textsuperscript{1}Amherst College
}
{
This study investigates how the direction of water flow influences habituation in oppositely oriented hair cells in the lateral line of larval zebrafish by pairing genetically encoded calcium indicators with fluorescent microscopy to visualize calcium transients in hair cells. Examining directional selectivity in habituation can provide insights into zebrafish behavior, from predator avoidance to schooling, and further our understanding of how sensory systems adapt to repeated stimuli.
}
{
This study uses 4–7 days post-fertilization zebrafish larvae with MyoGCaMP8m, a genetically encoded calcium indicator, encoded in their hair cells. Prior to measuring calcium signals, the larvae are submerged in MS-222, a temporary anesthetic. The larvae are then immobilized with tungsten pins and injected with α-bungarotoxin to prevent motion artifacts during imaging. After immobilizing the larvae, calcium transients were visualized in lateral-line hair cells in response to fluid jet stimulation. Hair cells were stimulated with the fluid jet in both anterior-to-posterior and posterior-to-anterior directions, and fluorescence was imaged using fluorescent microscopy. Imaging data were acquired with Micro-Manager and Igor software, and analyzed using Fiji to quantify fluorescence over time.
}
{
The results of this study indicate that hair cells oriented in the A-to-P direction displayed lower peak fluorescence and a faster decline in calcium signal over time compared to P-to-A oriented cells. Further, when comparing peak fluorescence relative to the previous baseline, posterior-facing hair cells exhibited faster rates of decreased fluorescence, suggesting greater habituation compared to anterior-facing hair cells. These findings suggest orientation-specific differences in repetitive stimulation between posterior and anterior-facing hair cells.
}
{
This study provides preliminary evidence that hair cell orientation in zebrafish lateral-line neuromasts influences habituation. The different calcium responses between A-to-P and P-to-A hair cells suggest that orientation plays a role in how zebrafish respond to repetitive stimuli. These findings support further investigation into how sensory systems adapt to stimulus direction and frequency, how calcium signals differ with agonists, and whether the habituation asymmetry exists at the afferent neuron level.
}

\newpage
